\chapter{Introduction}
\label{ch:introduction}

\section{Background}
\label{sec:background}

Modern industrial production environments depend on standardised digital templates to
generate consistent product labels and communication materials across global supply
chains~\cite{zhong2019publaynet}. These templates define the spatial organisation of
textual and graphical elements, including product identifiers, barcodes, regulatory symbols,
material composition statements, and country-of-origin declarations. Maintaining structural
and content consistency across templates is critical for product traceability, regulatory
compliance, and operational efficiency.

In template-driven workflows, structured metadata—commonly encoded as
\acl{JSON}~(JSON)—is used to dynamically populate predefined label designs. While this
approach improves scalability and automation, it introduces strict requirements for layout
accuracy: each generated label must follow predefined positioning rules and visual constraints
to ensure that all elements remain readable, correctly aligned, and compliant with design
standards~\cite{li2020docbank}. Even a small positional shift—such as a text field
overlapping a barcode quiet zone, or an element migrating outside its designated region—can
cause a label to be rejected by the production system or by regulatory authorities.

Recent advances in deep learning have transformed document image analysis and opened
practical routes to automating such compliance tasks. Datasets such as
PubLayNet~\cite{zhong2019publaynet} and DocLayNet~\cite{pfitzmann2022doclaynet} have
enabled the training of neural networks capable of identifying structural regions within
complex document layouts. Frameworks such as LayoutParser~\cite{shen2021layoutparser}
make these models accessible to practitioners in industrial settings. Transformer-based
models such as LayoutLM~\cite{xu2020layoutlm} and LayoutLMv3~\cite{huang2022layoutlmv3}
further show that jointly encoding text, layout, and visual information yields strong
document-understanding performance. Despite these advances, most existing solutions focus
on general-purpose document understanding rather than deterministic template validation
against approved reference designs.

In this thesis, the problem is investigated within the context of IKEA of Sweden's label
production workflow. Labels are engineered in NiceLabel and saved as proprietary
\texttt{.nlbl} files. Before deployment, each label must be verified against an approved
baseline template by human quality assurance (QA) testers. As the product range and the
number of template variants grows, manual verification becomes a bottleneck that motivates
the development of an automated alternative.

\section{Problem Statement}
\label{sec:problem}

Despite increasing automation in industrial production workflows, ensuring the correctness
of generated templates remains a significant technical challenge. In many modern systems,
templates are dynamically generated using structured input data, where both visual layout
and content integrity must remain consistent with predefined standards—approved baseline
label templates, internal business rules, barcode and regulatory placement requirements,
and production-specific design constraints. In the context of this thesis, label outputs are
engineered in NiceLabel and reviewed against strict reference designs before deployment,
meaning that even small deviations in spacing, alignment, font scaling, element positioning,
or metadata placement may cause a label to be rejected.

A central difficulty lies in the complex dependencies between layout structure, metadata,
and visual elements. Changes to a single element—such as modifying text length or
repositioning a field—may unintentionally affect surrounding elements, resulting in
cascading layout inconsistencies. Existing document layout analysis systems have made
significant progress in detecting layout components, yet many are optimised for document
understanding rather than deterministic compliance checking against predefined reference
templates~\cite{pfitzmann2022doclaynet,shen2021layoutparser}.

Another challenge arises from the limited availability of labelled failure data. Incorrect
templates are corrected immediately in production rather than archived, making supervised
learning from historical negative examples impractical. Validation strategies that combine
structured metadata verification with visual comparison against approved references are
therefore more suitable than approaches that depend on large collections of labelled errors.

\section{Motivation}
\label{sec:motivation}

The motivation for this research originates from the increasing demand for automation in
industrial template validation workflows. In the current process, trained human testers
review each generated label against approved standards—a process estimated to take an
average of 8–12 minutes per label. At production scale, with multiple template revisions
per product launch, this represents a significant allocation of skilled human time and
introduces the risk of inconsistent judgements.

Automated validation systems provide repeatable, deterministic verification, reducing
variability introduced by human interpretation. By ensuring that templates are evaluated
against consistent criteria, organisations can achieve higher levels of quality assurance,
shorter validation cycles, and lower operational cost. The development of an automated
system also supports long-term scalability as the template catalogue grows.

\section{Research Questions}
\label{sec:rqs}

This thesis is guided by the following main research question and three sub-questions:

\begin{tcolorbox}[title=Main Research Question]
\textbf{RQ1:} How can an automated validation system be designed to verify layout
structure, metadata correctness, and visual consistency in standardised product label
templates?
\end{tcolorbox}

\begin{itemize}
  \item \textbf{RQ1.1} How can reference-based comparison against approved baseline
  templates be used to detect layout deviations in generated labels?

  \item \textbf{RQ1.2} How can rule-based validation of structured metadata be used to
  detect logical and content-related inconsistencies before or during label validation?

  \item \textbf{RQ1.3} How can detected validation errors be localised and reported in a
  way that supports manual correction of the label template?
\end{itemize}

\section{Scope and Limitations}
\label{sec:scope}

The scope of this thesis is focused on:
\begin{itemize}
  \item Approximately 20 standardised label types from IKEA of Sweden's current
  catalogue, each with multiple design versions.
  \item Rule-based validation of JSON metadata fields prior to rendering.
  \item Reference-based visual comparison of rendered label images against approved
  baseline PNGs, using SSIM, ORB feature matching, and CNN-based classifiers.
  \item Bounding-box localisation of detected violations to support human correction.
  \item A REST API exposing the pipeline for integration with existing production systems.
\end{itemize}

The following are explicitly out of scope:
\begin{itemize}
  \item Full physical dimension validation (pixel-to-millimetre calibration without a
  known reference object).
  \item Multi-language semantic validation of free-text label fields.
  \item Production-grade deployment with enterprise authentication, load balancing, and
  continuous operation.
  \item Large-scale deep learning training from scratch; pre-trained backbone networks
  are used with fine-tuning.
\end{itemize}

\section{Thesis Structure}
\label{sec:structure}

\textbf{Chapter 1 (Introduction)} presents the background, problem statement, and
research questions. \textbf{Chapter 2 (Literature Study)} reviews related work in document
layout analysis, visual comparison, and industrial quality inspection. \textbf{Chapter 3
(Methodology)} describes the system architecture, dataset, validation strategy,
ML algorithms, and evaluation method. \textbf{Chapter 4 (Results)} presents experimental
findings. \textbf{Chapter 5 (Remaining Work and Future Development)} outlines next steps.

\section{Thesis Authorship Statement}
\label{sec:authorship}

We confirm that this thesis represents our own original work. All concepts, system designs,
and analyses have been developed by the authors as part of the bachelor thesis project
conducted in collaboration with IKEA of Sweden. External sources are cited in accordance
with academic standards. Writing assistance tools were used solely for language refinement;
all technical and conceptual content is our own. Both authors participated equally in
system design, implementation, analysis, and documentation.
